\begin{example}[Repelling passage with an oblique exit]\label{ex:repelling}
To see the role of the margin $A-D$ in Lemma~\ref{lem:terminal-distortion},
consider a time-dependent scalar equation, or equivalently its planar
extension $\dot t=1$. Write $U=2+x$. Until $t=0$ take
$\epsilon x'=-x$ and launch $x=q\in(0,1)$ at $t=-A+\delta y$,
where $A>0$, $\delta\downarrow0$ and $\epsilon=o(\delta)$.
The reference is $x=0$. On $0\le t\le1$ use
$\epsilon x'=D x(1-x)$ with $D>0$. After $t=1$ take
$\epsilon U'=-1$ and stop on the oblique line $U=\alpha+\gamma(t-1)$,
with $0<\alpha<1$ and fixed $\gamma$. At $t=0$ the entering
distance is $d_y=q e^{(-A+\delta y)/\epsilon}$. The exact logistic
flow and the final straight flight give the terminal time gap
\begin{equation}\label{eq:logistic-gap}
 g_y=\frac{\epsilon}{1+\epsilon\gamma}\,
        \frac{e^{D/\epsilon}d_y}{1+(e^{D/\epsilon}-1)d_y}.
\end{equation}
On the last flight $U-\alpha-\gamma(t-1)$ straightens the section and
has uniformly negative fast velocity. On the repelling interval
$(f_U)_+\le D$. Thus the lemma applies with a strictly positive
margin for bounded $y$ when $D<A$.
More precisely, for fixed $A,D$ and $y$,
\[
 \lim\frac{\epsilon}{\delta}\log(g_y/g_0)=
 \begin{cases}
 y,&D<A,\\
 \min(y,0),&D=A,\\
 0,&D>A .
 \end{cases}
\]
Indeed the denominator in~\eqref{eq:logistic-gap} is
$1+q(1-e^{-D/\epsilon})e^{(D-A+\delta y)/\epsilon}$.
At equality its logarithm, multiplied by $\epsilon/\delta$, tends to
$\max(y,0)$. Above equality the denominators cancel the entire input
slope. At $D=A$ a two-sided window of launches includes inputs wider
than $q e^{-A/\epsilon}$ by the factor $e^{\delta y/\epsilon}$ for $y>0$.
On the smaller bundle $0\le d\le d_0=q e^{-A/\epsilon}$, differentiating
the passage in~\eqref{eq:logistic-gap} at $D=A$ gives
\[
  K'(0)/K'(d_0)=\big(1+q(1-e^{-A/\epsilon})\big)^2
       \longrightarrow(1+q)^2.
\]
Thus the oscillation of
$\epsilon\log K'$ at $D=A$ is
$2\epsilon\log(1+q(1-e^{-A/\epsilon}))$. This is of order $\epsilon$,
hence $o(\delta)$ on the smaller bundle. A nonnegative margin alone
cannot guarantee exponential distortion control.

For each sufficiently small $\epsilon>0$ a smooth planar field realizes
the same gaps. Replace the two rates by
smooth weights with disjoint time supports, equal to $1$ near the
launch for the contracting rate, with integrals $A$ on $[-A,0]$
and $D$ on $[0,1]$, respectively. After $t=1$ use fast velocity
$-\chi((t-1)/\epsilon)$, with $\chi$ smoothly increasing from zero
to one on $[0,1]$, constant outside. The common displacement during
this ramp is at most one. Since $\alpha<1$, the ramp stays a fixed
distance from the exit for small $\epsilon$. All hits occur on the
following constant-speed flight and still satisfy~\eqref{eq:logistic-gap}.
\end{example}
